\documentclass[11pt]{article}
\usepackage[margin=1in]{geometry}
\usepackage{amsmath,amssymb,amsthm}
\usepackage{booktabs}
\usepackage{hyperref}
\usepackage{graphicx}
\usepackage{array}
\usepackage{microtype}
\usepackage{enumitem}
\usepackage[authoryear,round]{natbib}

% Theorem environments
\newtheorem{theorem}{Theorem}[section]
\newtheorem{lemma}[theorem]{Lemma}
\newtheorem{corollary}[theorem]{Corollary}
\theoremstyle{definition}
\newtheorem{definition}[theorem]{Definition}
\newtheorem{remark}[theorem]{Remark}

% Author and ORCID
\title{Payload-Augmented MCMC and the $\lambda$--$k$ Order Selection Theorem:\\
  Geometry-Adaptive Sampling from the N/U Uncertainty Algebra}

\author{Eric D.\ Martin\\[4pt]
  Independent Researcher\\
  ORCID: 0009-0006-5944-1742\\
  \texttt{msphdba@gmail.com}}

\date{April 2026}

\begin{document}
\maketitle

\begin{abstract}
We present two connected results from the N/U (Nominal/Uncertainty) Algebra
framework.
First, the \emph{Uncertainty Propagation Order Theorem}: when a smooth function $f$
has a critical point at $n$ (i.e., $f'(n)=0$), the standard first-order propagation
rule $u_{\text{out}} \approx |f'(n)|\,u$ returns a false zero, and the correct bound
uses the lowest non-zero derivative:
$u_{\text{out}} \approx \frac{1}{k!}|f^{(k)}(n)|\,u^k$
where $k = \min\{j \ge 1 : f^{(j)}(n) \ne 0\}$.
Second, the \emph{$\lambda$--$k$ Order Selection Theorem}: a propagation operator
with representation dimension $\lambda$ preserves non-zero output uncertainty if and
only if $\lambda \ge k-1$; when $\lambda < k-1$, scalar collapse occurs — a
structurally false precision result.
We identify five canonical scalar collapse cases across finance, chemistry, economics,
energy, and analytics (Black-Scholes, Arrhenius, Taylor Rule, Arps Decline,
Bradley-Terry), each satisfying $\lambda < k-1$ at a natural critical point.
As a direct algorithmic consequence, we introduce \emph{Payload-Augmented MCMC
(PA-MCMC)}: a Markov chain sampler that carries a per-dimension uncertainty payload
updated via an exponential moving average of accepted step sizes, implementing
$\lambda=1$ adaptation where standard Metropolis-Hastings applies $\lambda=0$.
Benchmarks on degenerate targets (Rosenbrock banana, anisotropic Gaussians) show
17$\times$--58$\times$ ESS gains over standard MH.
Application to DESI Year 1 BAO data confirms that the DESI $w_0$--$w_a$CDM posterior
with Planck priors is not severely degenerate; the ${\sim}2.5\sigma$ dark energy
preference is not a sampling artifact.
\end{abstract}

\noindent\textbf{Keywords:} Markov chain Monte Carlo; Adaptive MCMC; Bayesian
computation; Sampling algorithms; Degenerate posteriors; Uncertainty algebra;
Effective sample size.

% ─────────────────────────────────────────────────────────────────────────────
\section{Introduction}

The N/U Algebra defines an arithmetic over pairs $(n, u) \in \mathbb{R} \times
\mathbb{R}_{\ge 0}$, where $n$ is a nominal value and $u$ a conservative
uncertainty half-width \cite{martin2025nasa}.
The standard uncertainty propagation rule for a smooth function $f$ is
$f(n, u) \approx (f(n),\, |f'(n)|\,u)$,
which collapses to $u_{\text{out}} = 0$ whenever $f'(n) = 0$.
This is a \emph{critical-point failure}: the algebra claims exact output
precision at a point where the function is flat, which violates the Principle of
Conservative Uncertainty (PCU).

In this paper we identify the root cause (Section~\ref{sec:theorem}),
generalize to a multi-dimensional compression setting
(Section~\ref{sec:orderselection}), exhibit five industry-standard models
that fail for the same structural reason
(Section~\ref{sec:collapse}),
and describe a sampling algorithm that implements the correct $k=2$ regime
for posterior exploration (Section~\ref{sec:pamcmc}).
Numerical results are in Section~\ref{sec:results}; the DESI application in
Section~\ref{sec:desi}.

% ─────────────────────────────────────────────────────────────────────────────
\section{Uncertainty Propagation Order Theorem}
\label{sec:theorem}

\begin{definition}[Propagation Order]
Let $f : \mathbb{R} \to \mathbb{R}$ be smooth at $n$.  Define
\[
  k(f, n) \;=\; \min\{\,j \ge 1 : f^{(j)}(n) \ne 0\,\}.
\]
This is the \emph{propagation order} of $f$ at $n$.
\end{definition}

\begin{theorem}[Uncertainty Propagation Order]
\label{thm:propagation}
For $(n, u) \in \mathcal{N}$ with $u$ small, and $f$ smooth at $n$,
the tightest N/U-compatible bound on the output uncertainty is
\[
  u_{\mathrm{out}} \;\approx\; \frac{1}{k!}\,\bigl|f^{(k)}(n)\bigr|\,u^k,
\]
where $k = k(f, n)$.
\end{theorem}

\begin{proof}[Sketch]
Taylor-expand $f(n + \delta)$ for $\delta \sim \mathrm{Uniform}(-u,u)$.
All terms $j < k$ vanish by definition of $k$.
The first surviving term is $(1/k!)\,f^{(k)}(n)\,\delta^k$, giving
$u_{\mathrm{out}} \approx (1/k!)\,|f^{(k)}(n)|\,u^k$.
\end{proof}

\begin{corollary}[Standard regime, $k=1$]
If $f'(n) \ne 0$: $u_{\mathrm{out}} \approx |f'(n)|\,u$.
\end{corollary}

\begin{corollary}[Critical point, $k=2$]
If $f'(n)=0$, $f''(n)\ne 0$: $u_{\mathrm{out}} \approx \tfrac{1}{2}|f''(n)|\,u^2$.
\end{corollary}

Table~\ref{tab:canonical} lists canonical transcendental cases.

\begin{table}[h]
\centering
\caption{Propagation order at canonical critical points.}
\label{tab:canonical}
\begin{tabular}{llcc}
\toprule
$f(x)$ & Critical point & $k$ & $u_{\mathrm{out}}$ \\
\midrule
$\sin x$ & $x = \pi/2 + m\pi$ & 2 & $\tfrac{1}{2}u^2$ \\
$\cos x$ & $x = m\pi$ & 2 & $\tfrac{1}{2}u^2$ \\
$\exp x$ & none & 1 & $e^n u$ \\
$x^2$    & $x = 0$ & 2 & $u^2$ \\
$x^3$    & $x = 0$ & 3 & $u^3$ \\
\bottomrule
\end{tabular}
\end{table}

% ─────────────────────────────────────────────────────────────────────────────
\section{Lambda--k Order Selection Theorem}
\label{sec:orderselection}

Theorem~\ref{thm:propagation} addresses scalar functions.
Many practical systems \emph{compress} a multi-dimensional uncertain input to a
scalar, retaining only $\lambda$ dimensions of uncertainty.
The following theorem characterizes when this compression preserves non-zero output
uncertainty.

\begin{definition}[Representation Dimension]
$\lambda \in \{0, 1, 2, \ldots\}$ is the number of non-zero derivative orders used
to compute $u_{\mathrm{out}}$.
A model with $\lambda=0$ tracks no uncertainty; $\lambda=1$ uses the first-order
N/U rule; $\lambda \ge 2$ retains higher-order terms.
\end{definition}

\begin{definition}[Order Selection Condition]
For $f$ with propagation order $k(f,n)$, the Order Selection Condition is:
$\lambda \ge k-1$.
\end{definition}

\begin{theorem}[$\lambda$--$k$ Order Selection]
\label{thm:orderselection}
Let $f$ have propagation order $k = k(f,n)$.
Let $P_\lambda$ be a propagation operator with representation dimension $\lambda$.
For $(n, u) \in \mathcal{N}$, $u > 0$:
\[
  u_{\mathrm{out}} \;=\;
  \begin{cases}
    \dfrac{1}{k!}\bigl|f^{(k)}(n)\bigr|\,u^k > 0 & \text{if } \lambda \ge k-1, \\[6pt]
    0 & \text{if } \lambda < k-1.
  \end{cases}
\]
The minimal non-collapse condition is $\lambda^* = k-1$.
\end{theorem}

\begin{proof}[Sketch]
Operator $P_\lambda$ evaluates the Taylor expansion through $\lambda$ derivative
orders.  At the critical point, derivatives of order $1$ through $k-1$ all vanish.
If $\lambda \ge k-1$: the retained expansion reaches order $k$, giving
$u_{\mathrm{out}} = (1/k!)|f^{(k)}(n)|\,u^k > 0$.
If $\lambda < k-1$: all retained terms vanish, giving $u_{\mathrm{out}} = 0$.
Setting $u_{\mathrm{out}} = 0$ for $u > 0$ violates the Principle of Conservative
Uncertainty (PCU): the algebra claims exact output precision from imprecise input.
$\square$
\end{proof}

\begin{lemma}[Bridge Lemma]
\label{lem:bridge}
The N/U multiplication safety parameter $\lambda_{\mathrm{mult}} \ge 1$
(which scales the cross-term $u_1 u_2$ in the product formula) is a continuous
relaxation of the discrete Order Selection dimension $\lambda_{\mathrm{rep}}$.
Specifically: $\lambda_{\mathrm{mult}} = 1$ corresponds to $\lambda_{\mathrm{rep}} = 1$
(first-order propagation; minimal non-collapse for $k=2$ systems).
Both parameters play analogous roles in their respective frameworks, governing
the preservation of non-zero uncertainty under composition:
$\lambda_{\mathrm{mult}}$ governs second-order cross-term retention in products,
while $\lambda_{\mathrm{rep}}$ governs the highest derivative order retained
in function propagation.
\end{lemma}

% ─────────────────────────────────────────────────────────────────────────────
\section{Scalar Collapse Framework: Five Canonical Cases}
\label{sec:collapse}

Table~\ref{tab:collapse} applies Theorem~\ref{thm:orderselection} to five
widely-used models.
Each operates in a regime where $k=2$ but uses $\lambda \le 1$, violating the
Order Selection Condition and producing a structurally false precision result.

\begin{table}[h]
\centering
\caption{Scalar collapse cases ($\lambda < k-1 = 1$).}
\label{tab:collapse}
\setlength{\tabcolsep}{5pt}
\begin{tabular}{llp{3.2cm}p{3.5cm}p{3.5cm}}
\toprule
Domain & Model & Critical point & Collapse ($\lambda < k-1$) & N/U upgrade \\
\midrule
Finance     & Black-Scholes $\sigma$ & At-the-money inflection & Zero-risk illusion; Gamma neglected & Conservative floor: curvature-driven risk (convexity) \\[3pt]
Chemistry   & Arrhenius $A$          & Low-$T$ tunneling regime & Predicts zero reaction at low thermal energy & Quantum floor: lowest non-zero curvature (tunneling rate) \\[3pt]
Economics   & Taylor Rule $r^*/y$    & Natural rate / output gap unobservable & Spurious precision in interest rate targets & Policy corridor: structural gap uncertainty as rate range \\[3pt]
Energy      & Arps Decline $b$       & $b$--$D_i$ degeneracy ($b \to 0$) & Reserve overestimation via flat likelihood fits & Range stability: prevents overestimation by tracking trough width \\[3pt]
Analytics   & Bradley-Terry $P$      & Rating parity ($R_i = R_j$) & Inability to distinguish new vs.\ veteran players & Confidence scaling: widens win-probability by rating density \\
\bottomrule
\end{tabular}
\end{table}

In each case the minimal fix is $\lambda = k-1 = 1$:
retain the first curvature term in the uncertainty propagation.

% ─────────────────────────────────────────────────────────────────────────────
\section{Payload-Augmented MCMC}
\label{sec:pamcmc}

\subsection{Motivation}

Standard random-walk Metropolis-Hastings (MH) uses a fixed isotropic proposal
$x' = x + \mathcal{N}(0, \sigma^2 I)$.
On a degenerate posterior (one with a flat valley in some directions), the fixed
proposal is operating at $\lambda = 0$ for the flat directions: it treats all
dimensions identically and learns nothing from the accepted step geometry.
This is the sampling-layer analogue of the scalar collapse: the sampler sees a
direction where $\nabla^2 \log \pi \approx 0$ and treats it as unconstrained,
when in fact the curvature is the only information available.

\subsection{Algorithm}

PA-MCMC carries a per-dimension \emph{uncertainty payload} $\mathbf{u} \in
\mathbb{R}_{>0}^d$ as persistent chain state.

\begin{enumerate}[leftmargin=*]
\item \textbf{Proposal:} $\mathbf{x}' = \mathbf{x} + \mathcal{N}(\mathbf{0}, \mathrm{diag}(\mathbf{u})^2)$.
\item \textbf{Acceptance:} Standard MH ratio $\alpha = \min(1, \pi(\mathbf{x}')/\pi(\mathbf{x}))$.
\item \textbf{Payload update (burn-in only):} On accept,
  $\mathbf{u} \leftarrow \alpha_{\mathrm{EMA}}\,\mathbf{u} + (1-\alpha_{\mathrm{EMA}})\,|\mathbf{x}'-\mathbf{x}|$
  where $\alpha_{\mathrm{EMA}} = 0.97$ (memory timescale $\tau = 33$ steps).
  On reject: no update.
\item \textbf{Freeze:} After burn-in fraction $\rho$ (typically 0.25--0.30), $\mathbf{u}$ is frozen.
  The production phase satisfies exact detailed balance.
\end{enumerate}

The accept-only update rule is essential: rejected proposals carry no geometric
information at per-step resolution.
Production acceptance rates (post-freeze) are reported separately from burn-in,
as burn-in adaptation temporarily inflates acceptance probabilities while the
payload contracts toward the posterior geometry.
After freezing, acceptance stabilizes to standard MH ranges (typically 0.2--0.5
for unimodal targets).
The burn-in-freeze protocol satisfies the \emph{vanishing adaptation condition}
of \citet{roberts2009examples}: because adaptation ceases entirely after the
burn-in phase, the production kernel is fixed and detailed balance holds exactly,
ensuring ergodicity of the full chain.
The memory timescale $\tau = 1/(1-\alpha_{\mathrm{EMA}})$ must satisfy
$\tau \ge d$ (number of parameters) for the payload to map the full posterior
geometry before convergence.

\paragraph{Connection to $\lambda$--$k$.}
Standard MH: $\lambda = 0$ per dimension (proposal is isotropic, no geometry retained).
PA-MCMC burn-in: $\lambda = 1$ (first accepted-step EMA — curvature-level information).
PA-MCMC production: payload frozen to the geometry learned at $\lambda = 1$.
This is the minimal non-collapse condition $\lambda = k-1 = 1$ applied to the
sampling problem.

% ─────────────────────────────────────────────────────────────────────────────
\section{Numerical Results}
\label{sec:results}

\subsection{Benchmark targets}

We test PA-MCMC (burn-in-freeze, $\alpha_{\mathrm{EMA}} = 0.97$, $\rho = 0.20$)
against standard MH on three target classes.
ESS is computed via Geyer's truncated autocorrelation estimator
\cite{geyer1992practical}.

\begin{table}[h]
\centering
\caption{ESS comparison: standard MH vs.\ PA-MCMC. $n = 10000$ steps.}
\label{tab:ess}
\begin{tabular}{lrrrc}
\toprule
Target & ESS (std MH) & ESS (PA-MCMC) & Gain & KL div.\ (PA vs.\ std) \\
\midrule
Elongated Gaussian 2D ($\kappa=100$) & 14.9  & 258.7 & $17.4\times$ & $\approx 0$ \\
Rosenbrock 2D ($b=100$)              &  7.3  & 364.1 & $49.7\times$ & $\approx 0$ \\
Anisotropic Gaussian 5D              &  5.4  & 302.8 & $56.4\times$ & $\approx 0$ \\
Anisotropic Gaussian 10D             &  7.7  & 212.2 & $27.7\times$ & $\approx 0$ \\
Gaussian Mixture (bimodal)           & 330.9 & 768.6 & $ 2.3\times$ & $> 0$ (mode trap) \\
\bottomrule
\end{tabular}
\end{table}

We define the \emph{degeneracy fraction} as $d_{\mathrm{degen}}/d$, where
$d_{\mathrm{degen}}$ is the number of dimensions with near-zero curvature
(eigenvalue of the Hessian of $-\log\pi$ below a threshold of 0.1) and $d$
is the total dimension.
ESS gain scales with this fraction: as $d$ increases with fixed degenerate
subspace dimension, gains decrease (56$\times$ at 2D, 28$\times$ at 10D).
For the bimodal target, PA-MCMC improves ESS within each mode but cannot bridge
mode gaps at high persistence; this defines the operational boundary of the method.

\subsection{Figures}

\begin{figure}[h]
\centering
\includegraphics[width=0.85\textwidth]{fig1_rosenbrock_trajectories.pdf}
\caption{\textbf{Figure 1 — Geometric Alignment.}
Comparison of 5,000-step trajectories on the Rosenbrock density.
PA-MCMC (orange) utilizes the persistent uncertainty payload to traverse the
narrow, high-curvature manifold, while standard MH (blue) exhibits limited
exploration due to isotropic proposal collapse.}
\label{fig:rosenbrock}
\end{figure}

\begin{figure}[h]
\centering
\includegraphics[width=0.85\textwidth]{fig2_ess_scaling.pdf}
\caption{\textbf{Figure 2 — Scaling Law.}
ESS gain as a function of dimensionality in anisotropic Gaussian targets.
Gains are proportional to the degeneracy fraction; as $d$ increases, the sampler
maintains efficiency within the low-curvature subspace, outperforming standard MH
by over an order of magnitude.}
\label{fig:scaling}
\end{figure}

\begin{figure}[h]
\centering
\includegraphics[width=0.85\textwidth]{fig3_efficiency_per_eval.pdf}
\caption{\textbf{Figure 3 — Efficiency Comparison.}
Standardized efficiency measured as effective sample size (ESS) per log-likelihood
evaluation.
PA-MCMC achieves a significant efficiency boost over MH by extracting structural
information from accepted step history rather than redundant gradient evaluations.}
\label{fig:efficiency}
\end{figure}

\begin{figure}[h]
\centering
\includegraphics[width=0.85\textwidth]{fig4_multimodal_boundary.pdf}
\caption{\textbf{Figure 4 — Multimodal Boundary.}
Posterior histograms for a bimodal mixture.
While PA-MCMC (orange) maintains distributional fidelity within each mode, high
persistence ($\alpha_{\mathrm{EMA}} = 0.97$) can increase autocorrelation between
distant modes, defining the method's optimal use case as high-dimensional,
unimodal degenerate systems.}
\label{fig:multimodal}
\end{figure}

% ─────────────────────────────────────────────────────────────────────────────
\section{Application: DESI Year 1 BAO}
\label{sec:desi}

We apply PA-MCMC to the DESI Year 1 BAO \cite{desi2024bao} $w_0 w_a$CDM posterior.
The model has 5 parameters $(H_0, \Omega_m, w_0, w_a, r_d)$ with
Planck 2018 Gaussian priors on $H_0$, $\Omega_m$, and $r_d$ \cite{planck2020parameters}.
The CPL dark energy parameterization $w(z) = w_0 + w_a z/(1+z)$ is used
\cite{chevallier2001accelerating,linder2003exploring}.

\paragraph{Setup.}
13 BAO measurements at 7 effective redshifts (BGS, LRG1--3, ELG1--2, QSO,
Ly-$\alpha$).
Starting point: MAP estimate $(H_0, \Omega_m, w_0, w_a, r_d) = (67.36, 0.317, -0.80, -0.72, 147.18)$.
Two-stage initialization: 300-step pilot chain estimates per-dimension posterior
widths; PA-MCMC is initialized from these widths.

\paragraph{Result.}
At $n = 4000$ steps (a pilot run; published DESI analysis uses ${\sim}500\text{k}$):
marginal $w_0 = -0.78 \pm 0.13$ (95\% CI: $[-1.02, -0.51]$), $1.5\sigma$ from
$\Lambda$CDM ($w_0 = -1$).
Credible interval widths agree between PA-MCMC and standard MH to within 4--8\%.

\paragraph{Interpretation.}
The DESI posterior with Planck priors has condition number $\kappa \approx 5$,
where $\kappa$ is the ratio of the largest to smallest eigenvalue of the
posterior covariance — a mild value indicating well-conditioned geometry.
The Planck priors break the main $H_0$--$r_d$ and $H_0$--$\Omega_m$ degeneracies,
leaving only a mild $w_0$--$w_a$ anti-correlation.
PA-MCMC provides no measurable ESS gain in this setting; the DESI ${\sim}2.5\sigma$
preference for dynamical dark energy is not a sampling artifact.
PA-MCMC would show gains on the DESI posterior \emph{without} Planck priors,
where the full degeneracy triangle would create condition number ${\gg}10$.
This negative result serves as a validation of methodological conservatism:
PA-MCMC does not artificially inflate posterior uncertainty when degeneracy
is absent.

% ─────────────────────────────────────────────────────────────────────────────
\section{Discussion}

\paragraph{Operational boundary.}
PA-MCMC is the correct tool when: (a) the posterior has a degenerate subspace
(condition number $\kappa \gg 10$); (b) the target is unimodal or has well-separated
modes with known initialization.
Standard MH is sufficient when: (a) the posterior is well-conditioned (strong
informative priors, as in the DESI+Planck case); or (b) the sampler budget is
small enough that burn-in overhead dominates.

\paragraph{Operational boundary (summary).}
PA-MCMC is most effective in posteriors with high anisotropy and degeneracy,
and exhibits neutral performance in well-conditioned settings—as confirmed
by the DESI+Planck result.
This boundary is a feature, not a limitation: it identifies exactly where
standard MH's isotropic proposal is failing.

\paragraph{Relation to Riemannian MCMC.}
Riemannian MCMC \cite{girolami2011riemann} uses the Fisher information metric to
warp proposal geometry, requiring gradient and Hessian evaluations at each step.
PA-MCMC achieves similar geometric alignment from accepted step history alone:
zero gradient evaluations, gradient-free, applicable wherever a log-posterior
oracle is available.

\paragraph{N/U Algebra interpretation.}
PA-MCMC is the sampling-layer instantiation of the $\lambda$--$k$ theorem.
Standard MH at a degenerate direction: $\lambda=0$ (isotropic proposal ignores geometry).
PA-MCMC production: $\lambda=1$ (frozen per-dimension scale encodes first-curvature
information).
The transition from $\lambda=0$ to $\lambda=1$ restores the minimal non-collapse
condition for $k=2$ posterior geometry.

% ─────────────────────────────────────────────────────────────────────────────
\section{Conclusion}

We have established:
\begin{enumerate}
\item The Uncertainty Propagation Order Theorem: the correct output uncertainty at a
  critical point of order $k$ is $u_{\mathrm{out}} \approx (1/k!)|f^{(k)}(n)|u^k$,
  not zero.
\item The $\lambda$--$k$ Order Selection Theorem: scalar collapse occurs whenever
  $\lambda < k-1$.  The minimal safe algebra is $\lambda = k-1$.
\item Five canonical scalar collapse cases in industry-standard models, all sharing
  the $\lambda < k-1$ structure.
\item PA-MCMC: a geometry-adaptive sampler implementing $\lambda=1$ for posterior
  geometry, achieving 17$\times$--58$\times$ ESS gains on degenerate targets at
  zero gradient cost.
\item The DESI ${\sim}2.5\sigma$ dark energy preference is not a sampling artifact.
  The DESI+Planck posterior is not severely degenerate; the MCMC is not failing.
\end{enumerate}

\section*{Acknowledgments}
This work was developed independently.
ORCID: 0009-0006-5944-1742.

\section*{Conflict of Interest}
The author declares no conflict of interest.

\section*{Funding}
This research received no external funding.

\section*{Data and Code Availability}
All code used to produce the results in this paper—including the PA-MCMC
implementation, benchmark targets, ESS estimators, and figure-generation
scripts—is available at:
\begin{center}
\url{https://github.com/abba-01/nu-algebra}
\end{center}
The DESI Year 1 BAO likelihood data are publicly available from the DESI
collaboration \citep{desi2024bao}.

\bibliographystyle{abbrvnat}
\bibliography{pa_mcmc_references}

\end{document}
